%%%%%%%%%%%%%%%%%%%%%%%%%%%%%%%%%
%
%       EXERCÍCIO
%
%%%%%%%%%%%%%%%%%%%%%%%%%%%%%%%%%

\definevalorquestao

\begin{exercicioBanco}[\valorquestao]
Uma fábrica considera aceitável que o comprimento \(x\), em milímetros, de uma peça satisfaça simultaneamente:
\begin{itemize}
    \item distância máxima de \(0{,}8\) mm do valor \(10\);
    \item distância mínima de \(0{,}3\) mm do valor \(9\).
\end{itemize}
Determine todos os valores possíveis de \(x\).

% Define as alternativas
\newcommand{\alternativas}{%
    \begin{center}
        \begin{tabularx}{\textwidth}{XXX}
            (a) \([9,3\,;\,10,8]\). &
            (b) \([9,2\,;\,8,7]\cup[9,3\,;\,10,8]\). &
            (c) \([9,2\,;\,8,7)\cup(9,3\,;\,10,8]\). \\[5pt]
            (d) \([9,2\,;\,10,8]\). &
            (e) \((-\infty,8,7]\cup[9,3,\infty)\).
        \end{tabularx}
    \end{center}
}

% Define a resposta correta
\newcommand{\resposta}{A}

% Lógica condicional para exibição
\ifdefstring{\atividade}{lista}{%
    \alternativas
    \vspace{0.5em}
    
    \noindent\textbf{Resposta:} letra \textbf{\resposta}.
}{%
    \ifdefstring{\modo}{objetivo}{%
        \alternativas
    }{%
        % modo = discursiva → não mostra alternativas
    }
}
\end{exercicioBanco}

